\documentclass[10pt,twocolumn]{article}
\usepackage[T1]{fontenc}
\usepackage[utf8]{inputenc}
\usepackage{lmodern}
\usepackage[margin=1.8cm,top=2.4cm,bottom=2cm,columnsep=0.6cm]{geometry}
\usepackage{fancyhdr}
\usepackage{titlesec}
\usepackage{booktabs}
\usepackage{amsmath,amssymb,amsfonts}
\usepackage{microtype}
\usepackage{xcolor}
\usepackage{mdframed}
\usepackage{parskip}
\usepackage{enumitem}
\usepackage{hyperref}

\definecolor{rulered}{RGB}{180,30,30}
\definecolor{boxgray}{RGB}{245,245,245}
\definecolor{keywordgray}{RGB}{100,100,100}

\pagestyle{fancy}
\fancyhf{}
\fancyhead[L]{\textsc{\small Claude Mag}}
\fancyhead[C]{\small\textit{Machine Learning Research Digest}}
\fancyhead[R]{\small 2026-08-24}
\fancyfoot[C]{\small\thepage}
\renewcommand{\headrulewidth}{0.8pt}

\titleformat{\section}{\normalsize\bfseries\color{rulered}}{}{0pt}{}%
  [\vspace{-4pt}\color{rulered}\rule{\columnwidth}{0.4pt}]
\titlespacing{\section}{0pt}{8pt}{4pt}

\hypersetup{colorlinks=true,urlcolor=rulered,linkcolor=black}

\begin{document}
\setlength{\parindent}{0pt}
\setlength{\parskip}{4pt}

{\small\color{keywordgray}\textbf{Keywords:}
  LLM reinforcement learning \textbullet{}
  credit assignment \textbullet{}
  value functions \textbullet{}
  privileged information}\par\vspace{4pt}

{\LARGE\bfseries\raggedright
  Le Critique: Privileged Value Functions
  for LLM Reinforcement Learning}\par\vspace{4pt}

{\large\itshape\raggedright
  Critics with Access to Task-Privileged Context Provide
  Token-Level Advantages Without Contaminating the Policy,
  Beating GRPO with Fewer Rollouts}\par\vspace{6pt}

{\small\textbf{By} Siddarth Venkatraman, Matthieu Dinot,
  Laurence Aitchison}\par\vspace{2pt}

{\small\color{keywordgray}arXiv:2608.16739 \textbullet{} Preprint, August 2026}
\par\vspace{2pt}\noindent{\color{rulered}\rule{\columnwidth}{0.6pt}}\vspace{6pt}

Group-relative methods like GRPO reduce gradient variance by sampling
multiple rollouts per prompt, but they assign reward only at sequence
level and are bottlenecked by straggler rollouts that reduce throughput
and inflate off-policyness. Learned value functions solve both problems
in principle, but training an accurate critic from the policy's own
context is notoriously hard. \textit{Le Critique} introduces Privileged
Value Functions (PVF), which give the critic access to task-privileged
information (ground-truth answer, rubric, execution state) unavailable
to the policy at inference time, enabling precise token-level advantage
estimates from a single rollout. An adaptive interpolation baseline
(TETHER) gracefully degrades to group-relative methods when the critic
is unreliable.

\begin{mdframed}[backgroundcolor=boxgray,linecolor=rulered,linewidth=1pt,
    innerleftmargin=6pt,innerrightmargin=6pt,
    innertopmargin=6pt,innerbottommargin=6pt]
\textbf{\small AT A GLANCE}\par\vspace{4pt}
\begin{description}[leftmargin=0pt,labelsep=4pt,itemsep=2pt,parsep=0pt,
    font=\small\bfseries]
  \item[Problem:] GRPO provides only sequence-level credit and suffers
    from straggler bottlenecks; standard value functions are hard to
    train from the policy's own context.
  \item[Why it matters:] Token-level credit assignment could dramatically
    accelerate LLM RL by identifying which reasoning steps actually
    contribute to task success.
  \item[Method:] Privileged Value Functions give critics access to
    task-privileged context (e.g., the ground-truth) during training
    only; TETHER adaptively interpolates between PVF and group baselines
    depending on estimated critic reliability.
  \item[Result:] $+4.2$ accuracy points on GSM8K/MATH vs.\ GRPO at the
    same training budget; $2.3\times$ faster convergence than PPO.
  \item[Evidence:] Ablations confirm that privileged context is essential;
    PVF without privileged information underperforms GRPO, while TETHER
    provides robustness when critic quality varies.
\end{description}
\end{mdframed}

\section{The Big Picture}

Reinforcement learning for LLMs improves model behavior on tasks with
scalar reward: math problem solving, code generation, instruction
following. The central challenge is efficient credit assignment: which
tokens in a response contributed to earning the reward, and by how much?

Credit assignment is easy when reward is dense (every token gets a
signal) or when rollouts are very short (a one-step response). For
multi-step reasoning, answer construction, or long tool-use sequences,
credit is sparse and long-delayed. The community has largely converged
on two approaches, each with a critical flaw.

\textbf{GRPO} avoids a critic by constructing a group of $G$ rollouts
and using group statistics as a baseline. The flaw: every token in a
response receives the same advantage, regardless of whether it was
a key reasoning step or filler text.

\textbf{PPO with a critic} provides token-level advantages but requires
training a separate model $V_\phi$ that can accurately predict return
from the policy's token context. The flaw: the context available to
the critic (the tokens generated so far) is the same context available
to the policy --- and the policy's context does not contain enough
information to reliably predict long-horizon return.

Le Critique's insight: the critic does not need to be limited to what
the \emph{policy} can see.

\section{Why Existing Approaches Fall Short}

\textbf{Sequence-level baselines (GRPO, REINFORCE++):} Assigning
the same advantage to all tokens in a 200-token reasoning chain means
the model cannot distinguish a critical deduction that led to the
correct answer from a redundant step that happened to precede it.
On math reasoning tasks, early reasoning steps (problem setup, variable
definition) are causally most important, but they receive the same
credit as conclusion steps under group-relative methods.

\textbf{Standard value functions (PPO):} The policy's context at step
$t$ contains the problem statement and all tokens generated so far.
This context does not directly contain the ground-truth answer, the
rubric used for scoring, or whether the model is ``on track'' toward
the correct solution. Trained on this incomplete information, critics
tend to be poorly calibrated: they predict high value for
high-confidence-but-wrong responses and fail to distinguish correct
from incorrect intermediate reasoning.

\section{Core Mechanism}

\textbf{Privileged Value Functions (PVF):} The critic $V_\phi(s_t, p_t)$
receives both the standard policy context $s_t$ (the token sequence)
and a privileged context $p_t$ available in the training environment
but not to the policy during inference. For math reasoning, $p_t$
might include the ground-truth answer and intermediate solution steps;
for code generation, it might include the test case outputs.

The privileged context is concatenated before the model's context
window at the critic input level and is never passed through the
policy's forward pass. The policy gradient uses advantages derived
from the PVF:
\[
  \hat{A}_t = r_t + \gamma V_\phi(s_{t+1}, p_{t+1}) - V_\phi(s_t, p_t)
\]
The policy parameters $\theta$ are updated via:
\[
  \nabla_\theta J = \mathbb{E}_t[\nabla_\theta \log\pi_\theta(a_t|s_t) \cdot \hat{A}_t]
\]
Gradients through the critic are blocked: $\phi$ is updated
separately with mean squared TD error.

\textbf{TETHER:} Not all critics are equally reliable. TETHER
adaptively interpolates between the PVF baseline and the
group-relative baseline:
\[
  b_t = \hat{\rho} \cdot V_\phi(s_t, p_t)
      + (1-\hat{\rho}) \cdot \mathbb{E}_G[R(\mathcal{T})]
\]
where $\hat{\rho} \in [0,1]$ estimates critic reliability
from the held-out validation loss:
\[
  \hat{\rho} = \sigma(-\alpha \cdot \mathcal{L}_V^{\text{hold-out}})
\]
High $\mathcal{L}_V$ (unreliable critic) drives $\hat{\rho} \to 0$,
recovering group-relative training; low $\mathcal{L}_V$ (reliable
critic) drives $\hat{\rho} \to 1$, using token-level PVF advantages.

\section{Technical Formulation}

Let $\pi_\theta$ be the LLM policy and $V_\phi$ the PVF critic.
During a training step, one rollout $\mathcal{T} = (s_0, a_0, \ldots,
s_T, r_T)$ is generated from $\pi_\theta$. The full TETHER policy
gradient is:
\[
  \nabla_\theta J = \mathbb{E}_{\mathcal{T}}\!\left[\sum_{t=0}^T
    \nabla_\theta \log\pi_\theta(a_t|s_t)
    \cdot \hat{A}_t^{\text{TETHER}}\right]
\]
\[
  \hat{A}_t^{\text{TETHER}} = r_t + \gamma b_{t+1} - b_t
\]
The critic is updated independently with:
\[
  \mathcal{L}_V = \mathbb{E}_t\!\left[(V_\phi(s_t,p_t)
    - (r_t + \gamma V_\phi(s_{t+1},p_{t+1})))^2\right]
\]

\section{Learning or Inference Procedure}

\textbf{Training:} Privileged context $p_t$ is constructed from
ground-truth task data (available in supervised datasets) and is
passed to the critic at each step. The critic is updated every
training iteration; reliability $\hat{\rho}$ is recomputed from
a 10\% held-out portion of each batch.

\textbf{Inference:} The critic and its privileged context are
discarded entirely. The policy $\pi_\theta$ runs without any
modifications, and privileged information is never exposed during
deployment.

\section{What the Guarantee Says}

No formal guarantee is provided for PVF. The key theoretical
observation is that the privileged critic, having access to the
ground truth, has a much lower Bayes-optimal value estimation error
than a critic trained from the policy's context alone. This tightened
value estimation translates to lower-variance advantage estimates,
which in turn reduces the sample complexity of policy improvement.

\section{Experimental Findings}

\begin{itemize}[itemsep=2pt,parsep=0pt]
  \item \textbf{Math reasoning (GSM8K + MATH):} PVF with TETHER
    achieves $+4.2$ accuracy points vs.\ GRPO at the same training
    budget (number of rollouts and gradient steps).
  \item \textbf{Training efficiency:} TETHER reaches GRPO's final
    accuracy $2.3\times$ faster in training steps, using
    fewer rollouts per gradient step.
  \item \textbf{Straggler reduction:} Eliminating the rollout-group
    requirement reduces training wall-clock time by 18--24\% on
    8-GPU configurations where straggler variance is highest.
  \item \textbf{PVF vs.\ standard critic:} PVF outperforms a standard
    PPO critic by $+2.9$ accuracy points, confirming that privileged
    context is the source of the advantage, not just having a critic.
\end{itemize}

\section{Ablations and Interpretation}

\textbf{Privileged context content:} Ground-truth answer provides
most of the benefit; adding intermediate solution steps provides
marginal additional gain. This suggests that knowing the target
outcome is the most valuable privileged information.

\textbf{TETHER vs.\ fixed interpolation:} Fixed $\hat{\rho}=0.5$
performs 1.8 pp worse than adaptive TETHER, confirming that
critic reliability varies during training and adaptivity is needed.

\textbf{Rollout group size:} TETHER with $G=1$ (single rollout)
outperforms GRPO with $G=8$ by 1.2 accuracy points, demonstrating
that privileged credit assignment compensates for reduced group size.

\textbf{Offline application:} PVF-trained policies retain their
advantage over GRPO baselines on out-of-distribution test
prompts, suggesting that the improved credit assignment produces
a qualitatively better-learned policy, not just in-distribution
fitting.

\section*{Reference}

Siddarth Venkatraman, Matthieu Dinot, Laurence Aitchison.
\textbf{Le Critique: Privileged Value Functions for LLM
Reinforcement Learning.}
arXiv:2608.16739, August 2026.
\url{https://arxiv.org/abs/2608.16739}

\end{document}
